\documentclass[12pt]{article}

\usepackage[utf8]{inputenc}
\usepackage[T1]{fontenc}
\usepackage[spanish,mexico]{babel}
\usepackage{amsmath,amssymb}
\usepackage{geometry}
\usepackage{enumitem}
\usepackage{booktabs}
\usepackage{longtable}
\usepackage{array}
\usepackage{xcolor}
\usepackage{hyperref}

\geometry{letterpaper,margin=2.5cm}
\setlist[itemize]{leftmargin=1.5em}
\setlist[enumerate]{leftmargin=1.8em}

\title{\textbf{Cálculo Diferencial}\\[0.3cm]
\large Tecnológico Nacional de México}
\author{Carlos E Mart\'inez Rodr\'iguez}
\date{}

\begin{document}
\maketitle
\tableofcontents

\section*{Informaci\'on del Curso}

\begin{itemize}
    \item \textbf{Asignatura:} Cálculo Diferencial.
    \item \textbf{Clave:} ACF--0901.
    \item \textbf{Horas teoría--práctica--créditos:} 3--2--5.
\end{itemize}

\section*{Objetivo general del curso}

El curso tiene como propósito que el/la estudiante \textbf{plantee y resuelva problemas utilizando las definiciones de límite y derivada de funciones de una variable para la elaboración de modelos matemáticos aplicados}.


\section{Números reales}

\textbf{Objetivos:} Aplicar las propiedades de los números reales, las desigualdades de primer y segundo grado con una incógnita y las desigualdades con valor absoluto, representando sus soluciones de manera gráfica y analítica.

\subsection{Subtemas}

\begin{enumerate}
    \item Propiedades de los números reales.
    \item Ecuaciones y desigualdades en los números reales.
    \item Desigualdades lineales.
    \item Desigualdades cuadráticas.
    \item Desigualdades con valor absoluto.
    \item Interpretación geométrica de desigualdades del tipo $|x-a|<b$.
    \item Planteamiento y resolución de problemas mediante desigualdades.
    \item Distinción entre números racionales y números reales mediante el axioma del supremo.
\end{enumerate}

\textbf{Tiempo previsto:} La planeaci\'on considera \textbf{13 horas} paraa esta unidad: una hora para presentación y evaluación diagnóstica y 12 horas efectivas de clase.

\section{Funciones}

\textbf{Objetivos:}  Analizar la definición de función real e identificar distintos tipos de funciones y sus representaciones gráficas para plantear modelos.

\subsection{Subtemas}

\begin{enumerate}
    \item Concepto de función real.
    \item Evaluación de funciones.
    \item Representación gráfica de funciones.
    \item Identificación y distinción de diferentes tipos de funciones.
    \item Operaciones con funciones.
    \item Obtención de funciones a partir del planteamiento de problemas simples.
    \item Uso de funciones para la construcción de modelos elementales.
\end{enumerate}

\textbf{Tiempo previsto:} La planeaci\'on considera \textbf{19 horas} para esta unidad: 18 horas de clase efectiva y una hora de examen escrito.

\section{Límites y continuidad}

\textbf{Objetivos:} Utilizar la definición de límite de funciones para determinar analíticamente la continuidad de una función en un punto o en un intervalo, así como mostrar gráficamente diferentes tipos de discontinuidad.

\subsection{Subtemas}

\begin{enumerate}
    \item Concepto de límite de una función.
    \item Interpretación gráfica del límite.
    \item Aplicación de la definición de límite a funciones lineales y cuadráticas.
    \item Cálculo de límites de funciones elementales mediante teoremas.
    \item Uso del límite para determinar el comportamiento de una función.
    \item Concepto de continuidad en un punto y en un intervalo.
    \item Identificación de discontinuidades mediante la definición.
    \item Representación gráfica de diferentes tipos de discontinuidad.
\end{enumerate}

\textbf{Tiempo previsto:} La planeaci\'on considera \textbf{13 horas} para esta unidad: 12 horas de clase efectiva y una hora para examen escrito.

\section{Derivadas}

\textbf{Objetivos:} Utilizar la definición de derivada para el análisis de funciones y el cálculo de derivadas.

\subsection{Subtemas}

\begin{enumerate}
    \item Concepto de derivada.
    \item Cálculo de derivadas mediante la definición.
    \item Interpretación geométrica de la derivada.
    \item Interpretación física de la derivada.
    \item Propiedades y teoremas de la derivada.
    \item Cálculo de derivadas utilizando reglas y propiedades.
    \item Derivación de funciones compuestas.
    \item Derivación implícita.
\end{enumerate}

\textbf{Tiempo previsto:} La planeaci\'on contempla \textbf{17 horas} para esta unidad: 16 horas de clase efectiva y una hora para examen escrito.

\section{Aplicaciones de la derivada}

\textbf{Objetivos:} Aplicar la derivada en la solución de problemas de optimización y de variación de funciones, y utilizar diferenciales en problemas que requieren aproximaciones.

\subsection{Subtemas}

\begin{enumerate}
    \item Teorema del valor medio.
    \item Aplicación de la derivada a problemas.
    \item Análisis de la variación de una función.
    \item Análisis e interpretación de la gráfica de una función mediante derivadas.
    \item Problemas de optimización.
    \item Concepto de diferencial.
    \item Aplicación de diferenciales a problemas de aproximación.
\end{enumerate}

\textbf{Tiempo previsto:} La planeaci\'on considera \textbf{18 horas} para esta unidad: 16 horas de clase efectiva y dos horas destinadas a evaluaciones escritas, una correspondiente al tema de derivadas y otra al tema de aplicaciones de la derivada.

\section{Planeaci\'on}

\begin{center}
\begin{tabular}{>{\raggedright\arraybackslash}p{8cm}c}
\toprule
\textbf{Unidad} & \textbf{Horas previstas} \\
\midrule
1. Números reales & 13 \\
2. Funciones & 19 \\
3. Límites y continuidad & 13 \\
4. Derivadas & 17 \\
5. Aplicaciones de la derivada & 18 \\
\midrule
\textbf{Total} & \textbf{80} \\
\bottomrule
\end{tabular}
\end{center}

\section{Criterios generales de evaluación}

En cada unidad se establece, de manera general, la siguiente ponderación:

\begin{itemize}
    \item \textbf{70\%}: evaluación temática.
    \item \textbf{30\%}: actividades de aprendizaje a criterio del docente.
\end{itemize}

Las actividades de enseñanza y aprendizaje contemplan la participación activa en clase, trabajo extra-clase, resolución de ejercicios, atención de dudas y evaluación escrita.

\section{Bibliograf\'ia b\'asica}


\begin{itemize}

\item Larson, Hostetler, Edwards, Cálculo I, Pirámide, 2002, 7a. 
\item Hughes-Hallett, Gleason et al., Cálculo aplicado, CECSA, 2003, 1a.   
\item Leithold, El cálculo, Oxford, 2003, 7a. 
\item Thomas, Finney, Cálculo. Una variable, Adison Wesley Longman, 1999, 9na.
\item Stewart, Cálculo. Trascendentes tempranas, Thomson, 2002,  4ta. 
\item Swokowski, Cálculo  con Geometría Analítica, Iberoamérica, 1989, 2da. 
 
\end{itemize}

\end{document}
